\documentclass[aspectratio=169]{beamer}
%--- Configuración del Tema y Colores ---
\usetheme{Madrid}
\usecolortheme{whale}
\setbeamertemplate{navigation symbols}{}
\setbeamertemplate{caption}[numbered]
%--- Paquetes ---
\usepackage[utf8]{inputenc}
\usepackage[spanish]{babel}
\usepackage{graphicx}
\usepackage{booktabs}
\usepackage{tikz}
\usepackage{fontawesome5}
\usepackage{listings}
\usetikzlibrary{shapes.geometric, arrows, positioning, calc, fit}

%--- Configuración de Listings ---
\lstset{
    basicstyle=\ttfamily\scriptsize,
    breaklines=true,
    frame=single,
    backgroundcolor=\color{gray!10},
    keywordstyle=\color{blue},
    commentstyle=\color{green!50!black},
    stringstyle=\color{red!70!black},
}

%--- Metadatos del Documento ---
\title[PM2]{Programación Matemática 2}
\subtitle{Clase 7: Ingeniería de Software - Docker y Testing}
\author[Luis Alfredo Alvarado]{Prof. Luis Alfredo Alvarado Rodríguez}
\institute[ECFM - USAC]{
    Escuela de Ciencias Físicas y Matemáticas \\
    Universidad de San Carlos de Guatemala
}
\date{Primer Semestre, 2026}

\begin{document}

%--- Diapositiva de Título ---
\begin{frame}
    \titlepage
\end{frame}

%--- Diapositiva: Tabla de Contenidos ---
\begin{frame}{Agenda de Hoy}
    \tableofcontents
\end{frame}

%===========================================================
% SECCIÓN 1: INTRODUCCIÓN A DOCKER
%===========================================================
\section{¿Qué es Docker?}

\begin{frame}{El Problema: ``En mi máquina sí funciona''}
    \centering
    \begin{tikzpicture}[node distance=1cm]
        \tikzstyle{machine} = [rectangle, rounded corners, minimum width=2.5cm, minimum height=1.5cm, text centered, draw=black, font=\small]
        \tikzstyle{problem} = [rectangle, minimum width=3cm, minimum height=0.6cm, text centered, draw=black, fill=red!20, font=\scriptsize]
        
        \node (dev) [machine, fill=green!20] {Laptop\\Desarrollador};
        \node (test) [machine, right=2cm of dev, fill=yellow!20] {Servidor\\Testing};
        \node (prod) [machine, right=2cm of test, fill=blue!20] {Servidor\\Producción};
        
        \node (p1) [problem, below=0.8cm of dev] {Python 3.11\\Ubuntu 22.04};
        \node (p2) [problem, below=0.8cm of test] {Python 3.9\\CentOS 7};
        \node (p3) [problem, below=0.8cm of prod] {Python 3.10\\Debian 11};
    \end{tikzpicture}
    
    \vspace{0.5cm}
    \textbf{Problemas comunes:}
    \begin{itemize}
        \item Diferentes versiones de lenguajes y librerías
        \item Configuraciones de sistema operativo distintas
        \item Dependencias faltantes o incompatibles
        \item Variables de entorno diferentes
        \item ``Funciona en desarrollo, falla en producción''
    \end{itemize}
\end{frame}

\begin{frame}{Docker: La Solución}
    \begin{block}{Definición}
        \textbf{Docker} es una plataforma que permite empaquetar aplicaciones junto con todas sus dependencias en unidades estandarizadas llamadas \textbf{contenedores}.
    \end{block}
    
    \vspace{0.3cm}
    \textbf{Beneficios principales:}
    \begin{itemize}
        \item \textbf{Consistencia:} Mismo entorno en desarrollo, testing y producción.
        \item \textbf{Aislamiento:} Cada contenedor es independiente.
        \item \textbf{Portabilidad:} Si funciona en Docker, funciona en cualquier lugar.
        \item \textbf{Eficiencia:} Más ligero que máquinas virtuales.
        \item \textbf{Reproducibilidad:} El entorno se define en código (Dockerfile).
    \end{itemize}
    
    \vspace{0.3cm}
    \centering
    \faDocker \ \textit{``Build once, run anywhere''}
\end{frame}

\begin{frame}{Contenedores vs Máquinas Virtuales}
    \centering
    \begin{tikzpicture}[node distance=0.1cm]
        \tikzstyle{layer} = [rectangle, minimum width=4cm, minimum height=0.7cm, text centered, draw=black, font=\scriptsize]
        
        % VM Stack
        \node (vm-title) at (-3, 3) {\textbf{Máquina Virtual}};
        \node (vm-hw) [layer, fill=gray!30] at (-3, 0) {Infraestructura};
        \node (vm-hyp) [layer, fill=orange!30, above=0.1cm of vm-hw] {Hypervisor};
        \node (vm-os1) [layer, fill=blue!20, above=0.1cm of vm-hyp, minimum width=1.8cm, xshift=-1cm] {Guest OS};
        \node (vm-os2) [layer, fill=blue!20, above=0.1cm of vm-hyp, minimum width=1.8cm, xshift=1cm] {Guest OS};
        \node (vm-app1) [layer, fill=green!20, above=0.1cm of vm-os1, minimum width=1.8cm] {App A};
        \node (vm-app2) [layer, fill=green!20, above=0.1cm of vm-os2, minimum width=1.8cm] {App B};
        
        % Docker Stack
        \node (d-title) at (3, 3) {\textbf{Docker}};
        \node (d-hw) [layer, fill=gray!30] at (3, 0) {Infraestructura};
        \node (d-os) [layer, fill=blue!20, above=0.1cm of d-hw] {Sistema Operativo Host};
        \node (d-docker) [layer, fill=cyan!30, above=0.1cm of d-os] {Docker Engine};
        \node (d-app1) [layer, fill=green!20, above=0.1cm of d-docker, minimum width=1.2cm, xshift=-1.3cm] {App A};
        \node (d-app2) [layer, fill=green!20, above=0.1cm of d-docker, minimum width=1.2cm] {App B};
        \node (d-app3) [layer, fill=green!20, above=0.1cm of d-docker, minimum width=1.2cm, xshift=1.3cm] {App C};
    \end{tikzpicture}
    
    \vspace{0.3cm}
    \begin{columns}[t]
        \column{0.48\textwidth}
        \begin{block}{VM}
            \begin{itemize}
                \item OS completo por VM
                \item GB de tamaño
                \item Minutos para iniciar
            \end{itemize}
        \end{block}
        
        \column{0.48\textwidth}
        \begin{block}{Contenedor}
            \begin{itemize}
                \item Comparte kernel del host
                \item MB de tamaño
                \item Segundos para iniciar
            \end{itemize}
        \end{block}
    \end{columns}
\end{frame}

%===========================================================
% SECCIÓN 2: CONCEPTOS FUNDAMENTALES DE DOCKER
%===========================================================
\section{Conceptos Fundamentales de Docker}

\begin{frame}{Conceptos Clave}
    \begin{columns}[t]
        \column{0.48\textwidth}
        \begin{block}{Imagen (Image)}
            \begin{itemize}
                \item Plantilla de solo lectura
                \item Contiene código, librerías, dependencias
                \item Se construye desde un Dockerfile
                \item Puede basarse en otras imágenes
            \end{itemize}
        \end{block}
        
        \begin{block}{Contenedor (Container)}
            \begin{itemize}
                \item Instancia ejecutable de una imagen
                \item Aislado del sistema host
                \item Tiene su propio filesystem
                \item Puede crearse, iniciarse, detenerse, eliminarse
            \end{itemize}
        \end{block}
        
        \column{0.48\textwidth}
        \begin{block}{Dockerfile}
            \begin{itemize}
                \item Archivo de texto con instrucciones
                \item Define cómo construir la imagen
                \item Reproducible y versionable
            \end{itemize}
        \end{block}
        
        \begin{block}{Registry (Docker Hub)}
            \begin{itemize}
                \item Repositorio de imágenes
                \item Docker Hub: registry público
                \item Imágenes oficiales verificadas
                \item Similar a GitHub para código
            \end{itemize}
        \end{block}
    \end{columns}
\end{frame}

\begin{frame}{Flujo de Trabajo con Docker}
    \centering
    \begin{tikzpicture}[node distance=1.5cm, auto]
        \tikzstyle{element} = [rectangle, rounded corners, minimum width=2cm, minimum height=1cm, text centered, draw=black, font=\small]
        \tikzstyle{arrow} = [thick,->,>=stealth]
        
        \node (dockerfile) [element, fill=yellow!20] {Dockerfile};
        \node (image) [element, right=2cm of dockerfile, fill=blue!20] {Imagen};
        \node (container) [element, right=2cm of image, fill=green!20] {Contenedor};
        \node (registry) [element, above=1.5cm of image, fill=purple!20] {Registry\\(Docker Hub)};
        
        \draw [arrow] (dockerfile) -- node[above, font=\scriptsize] {\texttt{docker build}} (image);
        \draw [arrow] (image) -- node[above, font=\scriptsize] {\texttt{docker run}} (container);
        \draw [arrow] (image) -- node[left, font=\scriptsize] {\texttt{docker push}} (registry);
        \draw [arrow] (registry) -- node[right, font=\scriptsize] {\texttt{docker pull}} (image);
    \end{tikzpicture}
    
    \vspace{0.5cm}
    \begin{exampleblock}{Analogía}
        \textbf{Imagen} = Receta de cocina \quad | \quad \textbf{Contenedor} = Plato preparado
    \end{exampleblock}
\end{frame}

\begin{frame}[fragile]{Anatomía de un Dockerfile}
\begin{lstlisting}[language=bash, title=Dockerfile]
# Imagen base
FROM python:3.11-slim

# Directorio de trabajo dentro del contenedor
WORKDIR /app

# Copiar archivos de dependencias
COPY requirements.txt .

# Instalar dependencias
RUN pip install --no-cache-dir -r requirements.txt

# Copiar el resto del codigo
COPY . .

# Puerto que expone la aplicacion
EXPOSE 8000

# Comando por defecto al iniciar el contenedor
CMD ["python", "app.py"]
\end{lstlisting}
\end{frame}


%===========================================================
% SECCIÓN 4: DOCKER COMPOSE
%===========================================================
\section{Docker Compose}

\begin{frame}{¿Qué es Docker Compose?}
    \begin{block}{Definición}
        \textbf{Docker Compose} es una herramienta para definir y ejecutar aplicaciones Docker multi-contenedor usando un archivo YAML.
    \end{block}
    
    \vspace{0.3cm}
    \centering
    \begin{tikzpicture}[node distance=0.8cm]
        \tikzstyle{container} = [rectangle, rounded corners, minimum width=2cm, minimum height=1cm, text centered, draw=black, font=\small]
        \tikzstyle{arrow} = [thick,<->,>=stealth]
        
        \node (web) [container, fill=green!20] {Web\\(Flask)};
        \node (db) [container, right=2cm of web, fill=blue!20] {Database\\(PostgreSQL)};
        \node (cache) [container, right=2cm of db, fill=red!20] {Cache\\(Redis)};
        
        \draw [arrow] (web) -- (db);
        \draw [arrow] (db) -- (cache);
        
        \node [below=1cm of db, font=\small] {\texttt{docker-compose.yml}};
        
        % Bounding box
        \node[draw, dashed, fit=(web)(db)(cache), inner sep=0.5cm, label=above:Aplicación Multi-Contenedor] {};
    \end{tikzpicture}
    
    \vspace{0.3cm}
    \textbf{Ventajas:} Un solo comando para levantar toda la infraestructura.
\end{frame}

\begin{frame}[fragile]{Ejemplo de docker-compose.yml}
\begin{lstlisting}[title=docker-compose.yml]
version: '3.8'

services:
  web:
    build: .
    ports:
      - "5000:5000"
    volumes:
      - .:/app
    environment:
      - DATABASE_URL=postgresql://db:5432/mydb
    depends_on:
      - db

  db:
    image: postgres:15
    environment:
      - POSTGRES_DB=mydb
      - POSTGRES_PASSWORD=secret
    volumes:
      - postgres_data:/var/lib/postgresql/data

volumes:
  postgres_data:
\end{lstlisting}
\end{frame}



%===========================================================
% SECCIÓN 5: INTRODUCCIÓN A TESTING
%===========================================================
\section{Introducción a Testing}

\begin{frame}{¿Por Qué Hacer Testing?}
    \begin{columns}
        \column{0.5\textwidth}
        \textbf{Sin tests:}
        \begin{itemize}
            \item ``Creo que funciona...''
            \item Miedo a hacer cambios
            \item Bugs descubiertos en producción
            \item Regresiones frecuentes
            \item Debugging manual constante
        \end{itemize}
        
        \column{0.5\textwidth}
        \textbf{Con tests:}
        \begin{itemize}
            \item Confianza en el código
            \item Refactorización segura
            \item Documentación ejecutable
            \item Detección temprana de bugs
            \item Integración continua posible
        \end{itemize}
    \end{columns}
    
    \vspace{0.5cm}
    \begin{block}{Regla de Oro}
        El costo de arreglar un bug aumenta exponencialmente según la etapa en que se descubre: desarrollo $<$ testing $<$ producción.
    \end{block}
\end{frame}

\begin{frame}{Tipos de Tests}
    \begin{table}[]
        \centering
        \renewcommand{\arraystretch}{1.3}
        \scriptsize
        \begin{tabular}{p{2.5cm} p{4.5cm} p{4.5cm}}
            \toprule
            \textbf{Tipo} & \textbf{Propósito} & \textbf{Ejemplo} \\
            \midrule
            \textbf{Unit Test} & Probar una función aislada & \texttt{test\_suma()} verifica que \texttt{2+2=4} \\
            \textbf{Integration Test} & Probar componentes juntos & API conecta correctamente a la DB \\
            \textbf{End-to-End (E2E)} & Probar flujo completo del usuario & Usuario puede registrarse y hacer login \\
            \textbf{Regression Test} & Verificar que bugs arreglados no vuelvan & Bug \#123 sigue corregido \\
            \textbf{Performance Test} & Medir velocidad y recursos & La API responde en < 200ms \\
            \textbf{Smoke Test} & Verificar funcionalidad básica & El servidor inicia correctamente \\
            \bottomrule
        \end{tabular}
    \end{table}
\end{frame}

%===========================================================
% SECCIÓN 6: PYTEST
%===========================================================
\section{Testing con Pytest}

\begin{frame}{¿Por Qué Pytest?}
    \begin{columns}
        \column{0.5\textwidth}
        \textbf{Ventajas de Pytest:}
        \begin{itemize}
            \item Sintaxis simple y legible
            \item Descubrimiento automático de tests
            \item Fixtures potentes y reutilizables
            \item Excelentes mensajes de error
            \item Gran ecosistema de plugins
            \item Compatible con unittest
        \end{itemize}
        
        \column{0.5\textwidth}
        \textbf{Instalación:}
        \begin{block}{}
            \texttt{pip install pytest}
        \end{block}
        
        \vspace{0.3cm}
        \textbf{Ejecutar tests:}
        \begin{block}{}
            \texttt{pytest} \\
            \texttt{pytest -v} (verbose) \\
            \texttt{pytest test\_file.py}
        \end{block}
    \end{columns}
\end{frame}

\begin{frame}[fragile]{Escribiendo Tests con Pytest}
\begin{columns}
\column{0.48\textwidth}
\begin{lstlisting}[language=Python, title=calculator.py]
def suma(a, b):
    return a + b

def dividir(a, b):
    if b == 0:
        raise ValueError(
            "No dividir por cero"
        )
    return a / b
\end{lstlisting}

\column{0.48\textwidth}
\begin{lstlisting}[language=Python, title=test\_calculator.py]
import pytest
from calculator import suma, dividir

def test_suma_positivos():
    assert suma(2, 3) == 5

def test_suma_negativos():
    assert suma(-1, -1) == -2

def test_dividir_normal():
    assert dividir(10, 2) == 5

def test_dividir_por_cero():
    with pytest.raises(ValueError):
        dividir(10, 0)
\end{lstlisting}
\end{columns}
\end{frame}



%===========================================================
% SECCIÓN 8: CI/CD
%===========================================================
\section{Integración y Despliegue Continuos}

\begin{frame}{CI/CD: Automatización del Flujo}
    \centering
    \begin{tikzpicture}[node distance=1cm, auto]
        \tikzstyle{step} = [rectangle, rounded corners, minimum width=1.8cm, minimum height=0.8cm, text centered, draw=black, fill=blue!20, font=\scriptsize]
        \tikzstyle{arrow} = [thick,->,>=stealth]
        
        \node (code) [step] {Código};
        \node (build) [step, right=0.6cm of code] {Build};
        \node (test) [step, right=0.6cm of build, fill=yellow!20] {Test};
        \node (review) [step, right=0.6cm of test] {Review};
        \node (staging) [step, right=0.6cm of review, fill=orange!20] {Staging};
        \node (prod) [step, right=0.6cm of staging, fill=green!20] {Production};
        
        \draw [arrow] (code) -- (build);
        \draw [arrow] (build) -- (test);
        \draw [arrow] (test) -- (review);
        \draw [arrow] (review) -- (staging);
        \draw [arrow] (staging) -- (prod);
        
        % CI/CD labels
        \draw [decorate, decoration={brace, amplitude=5pt, raise=3pt}] 
            (code.north west) -- (test.north east) node[midway, above=8pt] {\small CI};
        \draw [decorate, decoration={brace, amplitude=5pt, raise=3pt}] 
            (staging.north west) -- (prod.north east) node[midway, above=8pt] {\small CD};
    \end{tikzpicture}
    
    \vspace{0.5cm}
    \begin{columns}[t]
        \column{0.48\textwidth}
        \begin{block}{Continuous Integration (CI)}
            \begin{itemize}
                \item Merge frecuente a main
                \item Build automático
                \item Tests automáticos
                \item Feedback rápido
            \end{itemize}
        \end{block}
        
        \column{0.48\textwidth}
        \begin{block}{Continuous Deployment (CD)}
            \begin{itemize}
                \item Deploy automático a staging
                \item Deploy a producción (manual o auto)
                \item Rollback automático si falla
            \end{itemize}
        \end{block}
    \end{columns}
\end{frame}

\begin{frame}[fragile]{GitHub Actions: Ejemplo Básico}
\begin{lstlisting}[title=.github/workflows/test.yml]
name: Tests

on:
  push:
    branches: [main]
  pull_request:
    branches: [main]

jobs:
  test:
    runs-on: ubuntu-latest
    steps:
      - uses: actions/checkout@v4
      - name: Set up Python
        uses: actions/setup-python@v5
        with:
          python-version: '3.11'
      - name: Install dependencies
        run: pip install -r requirements.txt
      - name: Run tests
        run: pytest -v
\end{lstlisting}
\end{frame}

%===========================================================
% SECCIÓN 10: EJERCICIOS
%===========================================================
\section{Ejercicios Propuestos}

\begin{frame}{Ejercicios}
    \begin{enumerate}
        \setlength\itemsep{0.7em}
        \item \textbf{Docker Básico:} Crear un Dockerfile para una aplicación Python simple. Construir la imagen y ejecutar el contenedor.
        
        \item \textbf{Docker Compose:} Crear un \texttt{docker-compose.yml} que levante una aplicación web con una base de datos PostgreSQL.
        
        \item \textbf{Testing:} Escribir al menos 5 tests unitarios con pytest para un módulo de funciones matemáticas (factorial, fibonacci, es\_primo).
        
        \item \textbf{TDD:} Implementar una función \texttt{validar\_email()} usando TDD. Primero escribir los tests, luego la implementación.
        
        \item \textbf{(Bonus)} Configurar GitHub Actions para ejecutar los tests automáticamente en cada push.
    \end{enumerate}
\end{frame}

\begin{frame}{Recursos Adicionales}
    \begin{columns}[t]
        \column{0.48\textwidth}
        \textbf{Docker:}
        \begin{itemize}
            \item \faDocker \ \texttt{docs.docker.com}
            \item Docker Hub: imágenes oficiales
            \item Play with Docker (práctica online)
            \item ``Docker Deep Dive'' - Nigel Poulton
        \end{itemize}
        
        \vspace{0.3cm}
        \textbf{Testing:}
        \begin{itemize}
            \item \texttt{docs.pytest.org}
            \item ``Python Testing with pytest'' - B. Okken
            \item Real Python tutorials
        \end{itemize}
        
        \column{0.48\textwidth}
        \textbf{CI/CD:}
        \begin{itemize}
            \item GitHub Actions docs
            \item GitLab CI/CD
            \item CircleCI, Jenkins
        \end{itemize}
        
        \vspace{0.3cm}
        \textbf{Con IA:}
        \begin{itemize}
            \item Cursor puede generar Dockerfiles
            \item Copilot sugiere tests automáticamente
            \item Claude puede explicar errores de Docker
        \end{itemize}
    \end{columns}
\end{frame}

%--- Diapositiva Final ---
\begin{frame}
    \centering
    \Huge \textbf{¿Preguntas?}
    
    \vspace{0.8cm}
    \large
    \textbf{Próxima clase:} Hardware para IA: CPU vs GPU vs TPU vs LPU
    
    \vspace{0.5cm}
    \normalsize
    Prof. Luis Alfredo Alvarado Rodríguez\\
    \small Escuela de Ciencias Físicas y Matemáticas
    
    \vspace{0.5cm}
    \footnotesize
    \textit{``If it works on your machine, then ship your machine.''}\\
    \textit{— O mejor aún, usa Docker.}
\end{frame}

\end{document}